\subsection{An explicit norm representative that cannot be continued}

\begin{statement}{A.98}{Proved}
{A non-surjective completed norm-fibre transition at the rational point three}
Let $\alpha=3^{1/3}>0$, $\gamma=3^{1/9}>0$,
$K=\Q(\alpha)$ and $E=\Q(\gamma)$, with the actual inclusion
$K\hookrightarrow E$ given by $\alpha=\gamma^3$.
Write $\widehat{F^\times}_3=\varprojlim_r F^\times/(F^\times)^{3^r}$
for each of these fields, and define
\[
 \mathcal F_F(-2)=
 \{z\in\widehat{F^\times}_3:\mathrm N_{F/\Q}(z)=\kappa_3(-2)\}.
 \tag{A.98.1}
\]
The continuous norm map
\[
 \mathrm N_{E/K}:\mathcal F_E(-2)\longrightarrow\mathcal F_K(-2)
 \tag{A.98.2}
\]
is not surjective, although both fibres are nonempty.  An explicit
element outside its image is the Kummer class of
\[
 \begin{aligned}
 u&=\frac{(4-\alpha)^3}{61},&
 \beta&=(1-\alpha)u
       =\frac{25-109\alpha+60\alpha^2}{61},\\
 \mathrm N_{K/\Q}(u)&=1,&
 \mathrm N_{K/\Q}(\beta)&=-2.
 \end{aligned}
 \tag{A.98.3}
\]

More generally, let $Q,R\in\Q^\times$ have zero valuation at
$61$, and let $e$ be a positive integer prime to three.
For the actual cubic extension
\[
 E_Q=K(\gamma_Q),\qquad \gamma_Q^3=Q^e\alpha,
 \tag{A.98.4}
\]
the completed element $\kappa_3(\beta R^e)$ is not in
$\mathrm N_{E_Q/K}(\widehat{E_Q^\times}_3)$.
These are exactly the norm-representative transitions associated
with the coefficient change $a=3\mapsto3Q^{3e}$ and
$b=-2\mapsto-2R^{3e}$.

Finally an ordinary local completed pair at $61$ whose second
coordinate is congruent to $-2$ modulo cubes cannot have positive
first-coordinate valuation there.  Thus the prime in this example
cannot be converted into positive $A$-support while preserving
that second-coordinate cube class.
\end{statement}
\begin{proof}
All norm maps in the statement are the usual field norms on
rational elements, followed by their continuous maps on power
completions.  Valuations induce continuous maps to $\Z_3$.
Under localization at a prime, the norm valuation is the sum of
residue degrees times the valuations at the primes above it.
We use that formula on the actual local factors.

The polynomial $T^3-3$ has the three distinct roots
\[
                          4,\quad5,\quad52\pmod {61}.
 \tag{A.98.5}
\]
Its derivative is nonzero at all three roots.  Hensel lifting
therefore gives three roots $\alpha_0,\alpha_1,\alpha_2\in\Q_{61}$
with these respective reductions and the product decomposition
\[
 K\otimes_\Q\Q_{61}\xrightarrow{\ \sim\ }\Q_{61}^3,
       \qquad f(\alpha)\longmapsto
                   (f(\alpha_0),f(\alpha_1),f(\alpha_2)).
 \tag{A.98.6}
\]
Since $4^3-3=61$, the first factor of $4-\alpha$ has valuation
one.  One can read this directly from
$(4-\alpha_0)(16+4\alpha_0+\alpha_0^2)=61$:
the second factor is a unit.  The other two factors of
$4-\alpha$ are units.  All factors of $1-\alpha$ are units.
Consequently the three valuations of $\beta$ are exactly
\[
                              (2,-1,-1).
 \tag{A.98.7}
\]
They remain these elements of $\Z_3$ after completion.
The identities
$\mathrm N_{K/\Q}(4-\alpha)=4^3-3=61$ and
$\mathrm N_{K/\Q}(1-\alpha)=1-3=-2$ prove (A.98.3).

Consider (A.98.4).  Its degree over $K$ is three:
at the prime above three the normalized valuation of
$Q^e\alpha$ is $3e\,v_3(Q)+1$, which is not divisible by
three, so this element is not a cube in $K$.
The coefficient comparison in the statement is given by the
literal equalities
\[
          \gamma_Q^9=3Q^{3e},\qquad
          \mathrm N_{K/\Q}(\beta R^e)=-2R^{3e}.
\]
At the three factors in (A.98.6), the local algebra is
\[
                 \Q_{61}[Z]/(Z^3-Q^e\alpha_i).
 \tag{A.98.8}
\]
The field $\Q_{61}$ contains $\mu_3$, by Hensel lifting
the three distinct residue roots of unity.
Each radicand in (A.98.8) is a unit.  A unit with noncube
residue gives an unramified cubic field: its residue polynomial
is irreducible of degree three and has unit derivative.
The norm valuations of that factor therefore belong to
$3\Z_3$.  A unit with cube residue is a cube in $\Q_{61}$
by Hensel lifting and gives the split algebra instead.

Multiplication of $\beta$ by $R^e$ does not change (A.98.7).
Every one of those valuations is nonzero modulo three.
For $\kappa_3(\beta R^e)$ to be a completed norm, all three
radicands $Q^e\alpha_i$ would consequently have to be cubes
in $\Q_{61}$.
Their ratios would then also be cubes.  However,
\[
 \frac{\alpha_1}{\alpha_0}\equiv\frac54=47\pmod {61},
                  \qquad 47^3=1,
                  \qquad47^{20}=13\ne1\pmod {61}.
 \tag{A.98.9}
\]
In the group $\F_{61}^{\times}$ of order sixty, every cube
has twentieth power one.  Equation (A.98.9) contradicts
the required cube condition and proves the general assertion.
Taking $Q=R=1$ proves the failure of surjectivity in (A.98.2).

Both fibres in (A.98.2) are nevertheless nonempty:
$1-\gamma$ has norm $-2$ to $\Q$, and its norm to $K$
is $1-\alpha$.  More precisely, for every $n\ge1$ put
$\alpha_n=3^{1/3^n}$ and $K_n=\Q(\alpha_n)$.
The norm polynomial gives the exact compatible family
\[
 \mathrm N_{K_{n+1}/K_n}(1-\alpha_{n+1})=1-\alpha_n,
          \qquad\mathrm N_{K_n/\Q}(1-\alpha_n)=-2.
 \tag{A.98.10}
\]
Thus this example records the failure of extending the specified
representative, while the same rational coordinate pair has the
displayed compatible family.  No nonzero coherent-norm obstruction
for that pair is asserted.

To check the final support assertion, one computes
\[
                         (-2)^{20}=47\ne1\pmod {61}.
 \tag{A.98.11}
\]
Thus $-2$ is not a cube in $\Q_{61}$.
For an ordinary local pair $(x,1-x)$ with $v_{61}(x)>0$,
the unit $1-x$ is in $1+61\Z_{61}$.
The cube map is an automorphism of that group, by simple-root
Hensel lifting at one.  Its cube class is therefore trivial,
contradicting the prescribed class of $-2$.
This proves the final assertion using the actual local Kummer map.
\end{proof}

The example applies directly to a proposed recursive norm method
that first chooses an arbitrary element of each completed norm
fibre and then tries to preserve it under further radical levels.
Even at the rational point three the indicated transition does
not preserve every choice.  A.95 instead constructs finite-level
divisor representatives and retains their convergence condition;
the local calculation above does not supply or refute that
convergence for its infinitely supported pair.

The exact certificate
\path{code/norm_representative_certificate.py}, with its retained
output, checks the polynomial norms, three local branches,
valuation vector and every scalar cube-class combination over
$\F_{61}$ used in this proof.
